\documentclass[11pt,a4paper]{article}

% ===== XeLaTeX + Korean =====
\usepackage{kotex}
\usepackage{fontspec}
\setmainfont{Times New Roman}
\setmainhangulfont{AppleMyungjo}
\setsanshangulfont{Apple SD Gothic Neo}

% ===== Layout =====
\usepackage[a4paper,top=18mm,bottom=18mm,left=20mm,right=20mm]{geometry}
\usepackage{fancyhdr}
\setlength{\headheight}{24pt}
\usepackage{enumitem}
\usepackage{mdframed}
\usepackage{array}
\usepackage{booktabs}

\setlength{\parindent}{1em}
\linespread{1.18}

% ===== Header / Footer =====
\pagestyle{fancy}
\fancyhf{}
\renewcommand{\headrulewidth}{0.6pt}
\fancyhead[L]{\sffamily\bfseries 2026 고1 영어 변형 — 정답 및 해설 (지문 26)}
\fancyhead[R]{\sffamily George Bird Grinnell}
\renewcommand{\footrulewidth}{0pt}
\fancyfoot[C]{\framebox[16mm][c]{\thepage}}

% ===== helpers =====
\newcommand{\qno}[1]{\par\vspace{10pt}\noindent{\sffamily\bfseries #1.}\ }
\newmdenv[linewidth=0.6pt,roundcorner=3pt,innertopmargin=7pt,innerbottommargin=7pt,
          innerleftmargin=9pt,innerrightmargin=9pt,skipabove=5pt,skipbelow=5pt]{pbox}

\begin{document}

\begin{center}
{\fontsize{15}{17}\selectfont\sffamily\bfseries [지문 26] 정답 및 해설 — George Bird Grinnell}
\end{center}
\vspace{6pt}

% ===== 정답 표 =====
\begin{center}
\renewcommand{\arraystretch}{1.3}
\begin{tabular}{ccccc}
\toprule
문항 & 유형 & 정답 & 배점 \\
\midrule
1 & 어법 네모 & ① & 2점 \\
2 & 내용 불일치 & ② & 2점 \\
3 & 서술형 단답(한글) & 아래 모범답안 참조 & 2점 \\
4 & 영어 제목 & ① & 2점 \\
\bottomrule
\end{tabular}
\end{center}

\vspace{8pt}
\hrule
\vspace{8pt}

% ===================== Q1 해설 =====================
\qno{1}{\sffamily 어법 네모 — 정답: ①\ \ developed / which / continues}

\vspace{4pt}
\noindent\textbf{(A) developed} (○) vs.\ \textit{was developing} ($\times$)

\noindent 어린 시절 새에 대한 관심을 \textit{갖게 되었다}는 과거의 일반적 사실이므로 단순과거 \textbf{developed}가 적절하다.
\textit{was developing}은 과거진행형으로, 어떤 행동이 진행 중에 다른 일이 끊어진 상황에서 쓰이므로 문맥상 부자연스럽다.

\vspace{6pt}
\noindent\textbf{(B) which} (○) vs.\ \textit{that} ($\times$)

\noindent 선행사 \textit{his first scientific expedition} 뒤에 쉼표(,)가 있으므로 계속적 용법의 관계대명사가 와야 한다.
계속적 용법에는 \textbf{which}만 쓸 수 있으며, \textit{that}은 계속적 용법에 사용 불가하다.

\vspace{6pt}
\noindent\textbf{(C) continues} (○) vs.\ \textit{is continued} ($\times$)

\noindent 주어 \textit{his work}는 단수이며, ``보전주의자들에게 영감을 주다''는 능동의 의미이므로 능동형 단수동사 \textbf{continues}가 적절하다.
\textit{is continued}는 수동태로, ``계속되어지다''라는 의미가 되어 문맥상 어색하다.

\vspace{8pt}
\hrule
\vspace{8pt}

% ===================== Q2 해설 =====================
\qno{2}{\sffamily 내용 불일치 — 정답: ②}

\vspace{4pt}
\noindent\textbf{② (거짓)} He participated in a scientific expedition \textit{prior to} completing his studies at Yale University.

\noindent 지문: ``\textit{After graduating from Yale University, he joined his first scientific expedition}''
→ 예일 대학교를 \textbf{졸업한 후에} 탐사에 참가했다. 선지 ②는 탐사 참가가 졸업 \textit{이전}이라고 했으므로 사실과 반대다.

\vspace{6pt}
\noindent\textbf{① (참)} 오듀본 공원(Audubon Park) 인근에서 성장했다는 내용과 일치한다. Audubon은 유명한 조류학자·자연주의자의 이름이다.

\noindent\textbf{③ (참)} 1870년대 후반부터 \textit{Forest and Stream}에 야생동물 보호를 촉구하는 기사를 많이 썼다는 내용과 일치한다.

\noindent\textbf{④ (참)} 1910년에 설립된 글레이셔 국립공원 설립에 기여했다는 내용과 일치한다. (20세기 초 = early twentieth century)

\noindent\textbf{⑤ (참)} 1938년 사망 후에도 그의 업적이 보전주의자들에게 영감을 준다는 내용과 일치한다.

\vspace{8pt}
\hrule
\vspace{8pt}

% ===================== Q3 해설 =====================
\qno{3}{\sffamily 서술형 단답(한글) — 모범답안 및 채점 기준}

\vspace{4pt}
\noindent\textbf{모범답안} (아래 중 2가지 이상 작성 시 만점):
\begin{enumerate}[leftmargin=2em, itemsep=3pt]
  \item 잡지 \textit{Forest and Stream}에 야생동물 보호를 촉구하는 기사를 씀
  \item 보전 단체(오듀본 협회 등) 설립에 기여함
  \item 글레이셔 국립공원 설립을 도움
  \item 초기 야생동물 보호법 제정에 영향을 미침
\end{enumerate}

\vspace{4pt}
\noindent\textbf{채점 기준}:
\begin{itemize}[leftmargin=2em, itemsep=2pt]
  \item 위 4가지 중 2가지 이상 정확히 서술 → 2점 (만점)
  \item 1가지만 정확히 서술 → 1점
  \item 0가지 또는 본문과 무관한 내용 → 0점
\end{itemize}

\vspace{8pt}
\noindent{\sffamily\bfseries 4. 제목 - 정답: ①}
\par 지문은 Grinnell의 생애 전체를 예일대 탐사, 보전 글쓰기, 보전 단체 설립, 국립공원 및 법 제정 영향까지 연결해 설명한다. 따라서 초기 야생동물 보호 운동의 목소리로서 Grinnell을 포괄하는 ①이 정답이다.

\end{document}
